\subsection{极小 Cauchy 滤子}

一致空间 X 上 Cauchy 滤子组成的集的（关于包含关系的）极小元素，称为 X 上的极小 Cauchy 滤子。

命题 5 设 X 是一致空间。对 X 上每个 Cauchy 滤子 F，存在唯一的粗于 F 的极小 Cauchy 滤子 \(\mathfrak{F}_0\)。若 B 是 F 的基，S 是 X 的对称一致邻域的一个基本系，则诸集合 V(M)（M ∈ B, V ∈ S）构成 \(\mathfrak{F}_0\) 的一个基。

若 M 与 \(M'\) 属于 B，V 与 \(V'\) 属于 S，则存在集合 \(M'' \in B\)（相应地 \(V'' \in S\)），使 \(M'' \subset M \cap M'\)（相应地 \(V'' \subset V \cap V'\)）；于是 \(V''(M'') \subset V(M) \cap V'(M')\)，因而诸集合 \(V(M)\)（ \(M \in B\)，\(V \in S\) ）确实构成 X 上一个滤子 \(F_0\) 的基。进而，若 M 是 V-小集，

则 V(M) 是 V-小集；因此 \(\mathfrak{F}_0\) 是 Cauchy 滤子，且显然粗于 \(\mathfrak{F}\)。为完成证明，只需证明：若 \(\mathfrak{G}\) 是粗于 \(\mathfrak{F}\) 的 Cauchy 滤子，则 \(\mathfrak{G}\) 细于 \(\mathfrak{F}_0\)。对每个 M ∈ \(\mathfrak{B}\) 与每个 V ∈ \(\mathfrak{S}\)，存在一个 V-小集 N ∈ \(\mathfrak{G}\)；由于 N ∈ \(\mathfrak{F}\)，N 与 M 相交；于是 N ⊂ V(M)，从而 V(M) ∈ \(\mathfrak{G}\)。

推论 1 对每个 \(x \in \mathbf{X}\)，邻域滤子 \(\mathfrak{B}(x)\)（\(x\) 在 \(\mathbf{X}\) 中的）是极小 Cauchy 滤子。

在命题 5 中取 \(\mathfrak{F}\) 为 \(\mathbf{X}\) 中含 \(x\) 的一切子集组成的滤子，并取 \(\mathfrak{B}\) 由单个元素 \(\{x\}\) 组成。

推论 2 Cauchy 滤子 F 的每个聚点 x 都是 F 的极限点。

存在一个滤子 \(\mathfrak{G}\)，它既细于 \(\mathfrak{F}\) 又细于 \(\mathfrak{B}(x)\)（第一章 § 7 第 2 目命题 4）；由于 \(\mathfrak{F}\) 是 Cauchy 滤子，\(\mathfrak{G}\) 也是。若 \(\mathfrak{F}_0\) 是粗于 \(\mathfrak{F}\) 的唯一的极小 Cauchy 滤子，则 \(\mathfrak{F}_0\) 与 \(\mathfrak{B}(x)\) 都是粗于 \(\mathfrak{G}\) 的极小 Cauchy 滤子。因此 \(\mathfrak{F}_0 = \mathfrak{B}(x)\)，这表明 \(\mathfrak{F}\) 收敛于 \(x\)。

推论 3 凡粗于某个收敛于点 x 的滤子的 Cauchy 滤子，也收敛于 x。

这是推论 2 的一个推论。

推论 4 若 \(\mathfrak{F}\) 是极小 Cauchy 滤子，则 \(\mathfrak{F}\) 的每个集合都有非空的内部，且该内部也属于 \(\mathfrak{F}\)（换言之，\(\mathfrak{F}\) 具有一个由开集组成的基）。

设 V 是 X 的任一一致邻域；则存在开的一致邻域 \(U \subset V\)（§ 1 第 2 目命题 2 的推论 2）。对 X 的每个子集 M，\(U(M)\) 是开集且含于 \(V(M)\)；于是由命题 5 即得结论。

